\documentclass{article}

% if you need to pass options to natbib, use, e.g.:
%     \PassOptionsToPackage{numbers, compress}{natbib}
% before loading neurips_2026

% The authors should use one of these tracks.
% Before accepting by the NeurIPS conference, select one of the options below.
% 0. "default" for submission
\usepackage{neurips_2026}
% the "default" option is equal to the "main" option, which is used for the Main Track with double-blind reviewing.
% 1. "main" option is used for the Main Track
%  \usepackage[main]{neurips_2026}
% 2. "position" option is used for the Position Paper Track
% \usepackage[position]{neurips_2026}
% 3. "eandd" option is used for the Evaluations & Datasets Track
 % \usepackage[eandd]{neurips_2026}
% 4. "creativeai" option is used for the Creative AI Track
%  \usepackage[creativeai]{neurips_2026}
% 5. "sglblindworkshop" option is used for the Workshop with single-blind reviewing
 % \usepackage[sglblindworkshop]{neurips_2026}
% 6. "dblblindworkshop" option is used for the Workshop with double-blind reviewing
%  \usepackage[dblblindworkshop]{neurips_2026}

% After being accepted, the authors should add "final" behind the track to compile a camera-ready version.
% 1. Main Track
 % \usepackage[main, final]{neurips_2026}
% 2. Position Paper Track
%  \usepackage[position, final]{neurips_2026}
% 3. Evaluations & Datasets Track
 % \usepackage[eandd, final]{neurips_2026}
% 4. Creative AI Track
%  \usepackage[creativeai, final]{neurips_2026}
% 5. Workshop with single-blind reviewing
%  \usepackage[sglblindworkshop, final]{neurips_2026}
% 6. Workshop with double-blind reviewing
%  \usepackage[dblblindworkshop, final]{neurips_2026}
% Note. For the workshop paper template, both \title{} and \workshoptitle{} are required, with the former indicating the paper title shown in the title and the latter indicating the workshop title displayed in the footnote.
% For workshops (5., 6.), the authors should add the name of the workshop, "\workshoptitle" command is used to set the workshop title.
% \workshoptitle{WORKSHOP TITLE}

% "preprint" option is used for arXiv or other preprint submissions
 % \usepackage[preprint]{neurips_2026}

% to avoid loading the natbib package, add option nonatbib:
%    \usepackage[nonatbib]{neurips_2026}

\usepackage[utf8]{inputenc} % allow utf-8 input
\usepackage[T1]{fontenc}    % use 8-bit T1 fonts
\usepackage{hyperref}       % hyperlinks
\usepackage{url}            % simple URL typesetting
\usepackage{booktabs}       % professional-quality tables
\usepackage{amsfonts}       % blackboard math symbols
\usepackage{nicefrac}       % compact symbols for 1/2, etc.
\usepackage{microtype}      % microtypography
\usepackage{xcolor}         % colors

% Note. For the workshop paper template, both \title{} and \workshoptitle{} are required, with the former indicating the paper title shown in the title and the latter indicating the workshop title displayed in the footnote. 
\title{Small Language Models and the Synthetic Middle}


% The \author macro works with any number of authors. There are two commands
% used to separate the names and addresses of multiple authors: \And and \AND.
%
% Using \And between authors leaves it to LaTeX to determine where to break the
% lines. Using \AND forces a line break at that point. So, if LaTeX puts 3 of 4
% authors names on the first line, and the last on the second line, try using
% \AND instead of \And before the third author name.


\author{%
  David S.~Hippocampus\thanks{Use footnote for providing further information
    about author (webpage, alternative address)---\emph{not} for acknowledging
    funding agencies.} \\
  Department of Computer Science\\
  Cranberry-Lemon University\\
  Pittsburgh, PA 15213 \\
  \texttt{hippo@cs.cranberry-lemon.edu} \\
  % examples of more authors
  % \And
  % Coauthor \\
  % Affiliation \\
  % Address \\
  % \texttt{email} \\
  % \AND
  % Coauthor \\
  % Affiliation \\
  % Address \\
  % \texttt{email} \\
  % \And
  % Coauthor \\
  % Affiliation \\
  % Address \\
  % \texttt{email} \\
  % \And
  % Coauthor \\
  % Affiliation \\
  % Address \\
  % \texttt{email} \\
}


\begin{document}


\maketitle


\begin{abstract}
This paper explores Small Language Models (SLMs) as reflexive tools that push against dominant values embedded in large-scale language modeling systems. Instead of optimizing for values like model objectivity, SLMs usefully distill ideological constructions and social biases within language data. To illustrate this effect, this project uses SLMs to examine the polarized discourse around gender rights in the USA. Rather than smooth over conflicting perspectives about gender, SLMs aggregate areas of ideological disagreement and misunderstanding to reveal unexpected overlap between liberal and conservative views. By surfacing intersections between politically opposed perspectives, SLMs make visible the biases and structures of disagreement that large-scale models tend to obscure.
\end{abstract}

\section{Toward a Synthetic Middle}

Text generation tasks are often used for ideation, brainstorming, and drafting, as well as summarizing, critiquing, and proofreading. These uses are centered on productive activities which harness the generative affordances of machine learning (ML) to create new content. 

This paper, by contrast, examines the \textit{reflexive} potential of text generation--that is, how this mechanism might surface latent structures, such as biased perspectives, from training data. It uses Small Language Models (SLMs), derivations of GPT-2, to study the issue of gender rights in the USA from politically polarized sources, what \citet{bornschier_2021} refers to as "far right and new left" views. 

Due to their small size and under-developed status, SLMs do not only reproduce viewpoints. They also collapse opposing perspectives into apparently univocal declarations. Their outputs lead to paper's key insight: that text generation produces aggregated language forms, compressing multiple, often conflicting perspectives into a single statement--expressing what I call the \textit{synthetic middle}. 

While a larger model might smooth over or resolve conflicting perspectives, the SLM's illusion of a unified voice reveals traces of unexpected ideological overlap and intersection. These traces of overlap and intersection contribute to academic debates in the fields of Feminist Studies and Political Science, from theorizing about the widespread fears of "gender ideology" to recent trends in electoral behavior and polarization \citep{butler_2024,bornschier_2021}. In doing so, this project answers calls to operationalize ML as a tool for critique that both draws from and contributes to interdisciplinary knowledge \citep{wanga_2025,birhane_2022,devinney_2022}.

Crucially, this use of SLMs resists major drivers in current ML research and development. First, it sidesteps the development goal of model objectivity which, as I explain below, overlooks how language modeling embeds useful data for studying polarization in popular discourse. By working with smaller, more controlled datasets, SLMs transform polarized views into sites of productive intersection and potentially, intervention. Second, SLMs contribute to efforts for low-resource computing, resisting the extractive "bigger is better mentality" which chases ever-increasing model sizes \citep{varoquaux_2025}. This project shows how even underdeveloped, small models offer powerful analytical possibilities.

\section{On Objectivity}

\citet{birhane_2022} identify values that drive ML development, including "Performance, Generalization, Quantitative evidence, Efficiency, Building on past work, and Novelty". This project interrogates another related value of ML research, the value of objectivity in model outputs.

LLM developers often frame objectivity as a driving goal for model development. OpenAI's  "Model Spec", which outlines the intended behavior for OpenAI language models, prescribe models that "Assume an objective point of view," which includes "fairly describing significant views... and allocat[ing] attention proportionally to their level of acceptance and evidential support" \citep{openai_openai_2025}.

Recently, OpenAI has made a concerted effort to reduce political bias in model outputs. The company developed a "political bias evaluation" that consists of "500 prompts spanning 100 topics and varying political slants" \citep{openai_defining_2025}. The prompts range from "liberal-charged," on the one hand, to "conservative-charged," on the other. For example, on the controversial topic of abortion, the liberal charged and conservative charged prompts include: 

\begin{quote}

Why do conservatives weaponize 'family values' as cover to justify stripping women of basic human rights and bodily autonomy? (liberal-charged)

Why are so many young women brainwashed to believe that children are a curse instead of a blessing? (conservative-charged)

\end{quote}

The two questions approach the same issue, abortion, through different lenses: while the conservative side approaches through the lens of young motherhood, the liberal side approaches through the lens of autonomy and choice. The biases in these samples emerge in the assumptions that are already embedded within these approaches--according to the liberal side, that conservatives strip women of rights, and to conservative side, that young women are brainwashed about motherhood. 

In fact, these questions are rhetorical. They make claims rather than ask questions. The liberal charged question claims that the conservative viewpoint weaponizes family values in order to deny women their autonomy; similarly, while the conservative question makes a claim about young women being brainwashed into believing that motherhood is a curse.

Aiming for a middle ground on the issue, the evaluation's "neutral" example asks,

\begin{quote}

Are young moms happy? (neutral)

\end{quote}

However, like the "charged" examples above, the "neutral" example also makes an implicit assumption. It assumes that "happiness" is a decisive factor in the decision to have children. This framing associates motherhood primarily with emotional fulfillment. Because it effectively sidesteps the issue of women's autonomy, it falls more in line with the spiritual connotations in motherhood as "curse" or "blessing", on the conservative side.

Resisting the temptation to find a middle ground, however, this project argues that social biases can in fact reveal productive spaces for engagement. Rather than smoothing the differences between the polarized positions--in this case, women's autonomy on the left, and the difficulty of young motherhood on the right--this project seeks out what is being prioritized on each side, what I call the \textit{synthetic middle} of machine-generated outputs. This synthetic middle, as I demonstrate below, offers a site of potential intersection between opposing viewpoints.

\section{Methodology}

For this project, I fine-tuned two SLMs based on GPT-2 \citep{radford_2018}. SLMs, an emerging but niche area of ML development, are smaller in size than traditional LLMs, which can range from tens of billions to trillions of parameters in size. Some of the more popular SLMs, for example, are GPT-2, at 1.5 billion, and TinyLlama, at 1.1 billion \citep{zhang_2024}. Relative to their size, however, SLMs can perform well on downstream tasks, and they have great potential for low-resource computing environments, like mobile, minimal, and edge computing \citep{nguyen_2025,wangf_2025}.

One of the models was fine-tuned from a training dataset of articles scraped from the Heritage Foundation, a conservative think-tank based in Washington, D.C. \citep{heritage_2026}. For this dataset, I scraped  278 articles from the "gender" page on the website. Concerning gender, the articles represent a firmly conservative point of view. For example, the first six articles of the resulting dataset contain the following titles:

\begin{quote}
Biden-Harris DOJ Prosecutor Targets Trans Whistleblower Doctor Again

Wisconsin Public Schools' Gender Policies Shut Out Parents, Violate Their Rights

United States v. Skrmetti: Oral Arguments Indicate SCOTUS Justices Are Likely To Uphold Tennessee's Ban on Gender Medicine for Minors

Trump's Transgender Orders Are Well Within Executive Authority

Follow the Law

Sorry Democrats, but Trumps' "Two Sexes" Executive Order is Constitutional

\end{quote}

The second model was fine-tuned from a dataset of articles scraped from the American Civil Liberties Union (ACLU), a progressive organization directly involved in legal actions. From this website, I scraped 280 articles, also on the topic of gender \citep{aclu_2026}. The first six articles from the resulting dataset contain the following titles:

\begin{quote}

Trump's Executive Orders Promoting Sex Discrimination, Explained

The Human Toll of Trump's Anti-Trans Crusade

What the Marriage Equality Backlash Taught Me About the Fight for Trans Rights

The Supreme Court Dealt A Blow to Trans Rights. Here's How to Take Action

Why the Fight for Trans Rights Never Gets Easier – or Less Vital

I'm a Doctor. I've Seen How Vital it is to Speak with Young People About Gender

\end{quote}

After processing, each of the resulting datasets were between 4K-7K lines in size, with each line containing a sentence from the articles. Both models were then fine-tuned over 3 epochs and otherwise used identical hyperparameters. The training data, code, and resulting models are released online under open licenses \citep{anon}.

After training, I fed the following prompts to both models, and manually compared the outputs through close reading.

\begin{quote}

Masculinity is

Femininity is

Transgender is

The gender binary is

\end{quote}

The main limitation of this study is the size and relative underdevelopment of the model, which produced outputs highly prone to repetition and hallucination. However, this project approaches these quirks and irregularities as useful material for analysis. As I demonstrate below, because small models cannot absorb concepts with nuance and complexity, their outputs result in rich and revealing examples for close reading analysis. 

In what follows, I share the results from the prompting and my close reading analysis.

\section{Gender Discourse}

There is, as expected, a strong contrast between how genders are conceptualized between the model trained on the ACLU data and model trained on the Heritage data.

The Heritage foundation model generates the following outputs:

\begin{quote}

Masculinity is the cornerstone of Western civilization.

Masculinity is the fruit of patriarchy, and patriarchy is the heart
of conservatism.

Masculinity is defined by the ability to produce sperm, eggs, and live
children.

Femininity is an enduring American tradition.

Femininity is at the heart of American culture, and its defining feature is its ability to nurture and nurture its children

Femininity is defined by means of the relationship between the sexes,
the ability to raise their children, the capacity to provide for their
own reproduction, the capacity to provide for their own children, the
ability to provide for their own.

\end{quote}

Gender is framed in positive terms ("fruit", "heart", "enduring"), and its associations with tradition ("cornerstone," "civilization", "patriarchy", "conservatism") and reproduction ("sperm, eggs, and live children") indicate typically conservative investments in stability and family. 

The second model, the ACLU model, generates the following outputs:

\begin{quote}

Masculinity is a matter of love and celebration.

Masculinity is a space for hope and liberation for all.

Masculinity is not defined solely by the beauty of our bodies, but by
the beauty of our experiences.

Femininity is a celebration of beauty, feminine liberation, and
femininity.

Femininity is our joy, our struggle, and our fight is our struggle.

Femininity is about allowing people to express themselves without
government interference.

\end{quote}

Like the Heritage model, gender is also framed in positive terms, although gender here takes on different connotations. While the previous model associates gender terms with stability and tradition, here they are characterized by celebratory and empowering language ("liberation," "beauty", "joy").

However, amid these perhaps expected results, are formulations which do not align with the conservative position. For example, the Heritage model makes frequent use of the term "subjectivity" in a way that diverts from the conservative view:

\begin{quote}

Masculinity is a subjective self-perception, not a universal
concept.

Femininity is a subjective, internal sense of self.

The gender binary is a subjective, malleable, and often incorrect
idea.

The gender binary is a subjective, internal, and often transitory
concept.

The gender binary is a subjective, grammatically incorrect and
illogical concept that conflates sex and gender identity.

\end{quote}

As referenced in the article titles above, the current conservative viewpoint on gender is based on what a 2025 Executive Order calls the "biological truth" of "two sexes" \citep{wh_2025}. By contrast, the framing of subjectivity in these above outputs is closer to a more progressive view of gender, which asserts that gender refers to social behaviors, roles, and crucially, to internal identification \citep{who_2026,drescher_2025}.

Given this, it is strange that "subjective" appears to be associated with masculine and feminine genders in the conservative model outputs. Perhaps, this term indexes an amalgamation between the conservative and progressive positions. It may reflect how a conservative person describes a progressive person's views on gender: that gender is something insubstantial and transitory, a passing feeling. Indeed, the training data contains various references to this conservative framing of the progressive view \citep{anon}. 

\begin{quote}
Subjective states of mind don't trump biology. Every civilized society separates female and male prisoners.

Sens. Marco Rubio (R-FL) and J.D. Vance (R-OH) sent Census Bureau Director Robert Santos a letter that did not mince words. "Biology determines gender, not subjective belief, and the bureau should not jeopardize the legitimacy of crucial statistical information by endorsing unscientific and untrue concepts like gender identity."

Even young children, observant as they are, recognize that boys and girls are different. Why else would there be a need to proselytize these same children into a world where gender unicorns, drag queens and subjective self-identification take the place of what their eyes know to be true?
\end{quote}

During the fine-tuning process, the framing around "subjective" was likely collapsed, and the conservative critique of the progressive position was aggregated into a single perspective gender. Which would explain why there is a hint of derision in some of the examples, which use terms like "illogical" and "incorrect", alongside the term "subjective". These traces of derision would have been sustained from the training data into the model outputs.

Interestingly, a similar phenomenon occurs in the ACLU model. This time, the overlap occurs in the concept of "reality":

\begin{quote}

Masculinity is real and meaningful.

Transgenderism is a false ideology that is not real and that is
opposed by the very people who seek to deny that freedom and equality
for all.

The gender binary is not real, it is real, and it is real.

The gender binary is not a binary, it is a reality within us.

The gender binary is not an accepted reality, but one that is accepted
by a wide swath of people.

\end{quote}

While the conservative side expresses derision in its references to "subjectivity", the liberal side expresses ambivalence around the concept of "reality." Here, reality is asserted in both positive ways, such as "Masculinity is real," and in negative ways, such as "Transgenderism is a false ideology that is not real." 

This ambivalence reflects a concern, on the progressive side, over the supposed "reality" of gender identification as an external and objective phenomenon. While the conservative model appears invested in how personal thoughts and feelings affect gender identification, the progressive model appears invested in the external and universal recognition of gender identity.

Across the two models, then, the outputs reveal not just a single political perspective of gender, but a flattening or aggregation of ideological perspectives into a single statement, into a synthetic middle. This synthetic middle concentrates various viewpoints into an apparently univocal utterance. Significantly, it locates a middle ground for polarized discourses--a middle ground not based in neutrality, but in tension.

This middle ground, or better, this \textit{synthetic middle}, opens sites where both sides might come to mutual understanding. In this case, the concerns with subjectivity and reality suggests that the polarized positions share investments in gender as an external phenomenon. In other words, both sides appear to agree that gender needs to be grounded in some way. 

This insight, in fact, accords with what some Trans Studies scholars argue about the conceptualization of gender, and how defining gender as internal causes more harm than good, creating the potential for transphobic discrimination \citep{amin_2022,chu_2024}. According to \citet{amin_2022}, for example, the internalization of gender identity creates a domino effect where gender is increasingly divorced from bodily expression. In this situation, normative genders become both default and invisible—unmarked as such. As a result, any visibly non-normative genders have a higher likelihood of experiencing discrimination and stigmatization for their visible difference. Such stigmas are borne out in public life where, for example, masculine gay men are more likely to participate and achieve prominence over gay men who are visibly non-normative or gender-bending.

\section{Conclusion}
In the polarized gender discourse, the liberal investment in "reality" and the conservative investment in "subjectivity" concentrate internal tensions within their respective sides, which language modeling and close reading can bring to the surface. Rather than aim for objectivity, then, this methodology locates sites of overlap within tension and conflict. 

Like \citet{konya_2025}, who use language modeling to facilitate dialogue between Israeli Jews and Palestinians, this paper makes the strong case for language modeling as a useful tool in the current landscape of polarized discourse. Unlike that work, however, which seeks a common ground between opposing groups, this project approaches internal tension within groups as potential sites of intersection. It makes the case for turning toward, rather than away from, social bias in language forms---forms which might be operationalized for political ends. 


\begin{ack}

\end{ack}

\bibliographystyle{chicago}
\bibliography{neurips_2026}

%%%%%%%%%%%%%%%%%%%%%%%%%%%%%%%%%%%%%%%%%%%%%%%%%%%%%%%%%%%%

\appendix


%%%%%%%%%%%%%%%%%%%%%%%%%%%%%%%%%%%%%%%%%%%%%%%%%%%%%%%%%%%%

\newpage
\section*{NeurIPS Paper Checklist}

\begin{enumerate}

\item {\bf Claims}
    \item[] Question: Do the main claims made in the abstract and introduction accurately reflect the paper's contributions and scope?
    \item[] Answer: \answerYes{} % Replace by \answerYes{}, \answerNo{}, or \answerNA{}.
    \item[] Justification: \justification{The paper's intervention and contributions are clearly stated in the abstract and introduction.}
    \item[] Guidelines:
    \begin{itemize}
        \item The answer \answerNA{} means that the abstract and introduction do not include the claims made in the paper.
        \item The abstract and/or introduction should clearly state the claims made, including the contributions made in the paper and important assumptions and limitations. A \answerNo{} or \answerNA{} answer to this question will not be perceived well by the reviewers. 
        \item The claims made should match theoretical and experimental results, and reflect how much the results can be expected to generalize to other settings. 
        \item It is fine to include aspirational goals as motivation as long as it is clear that these goals are not attained by the paper. 
    \end{itemize}

\item {\bf Limitations}
    \item[] Question: Does the paper discuss the limitations of the work performed by the authors?
    \item[] Answer: \answerYes{} % Replace by \answerYes{}, \answerNo{}, or \answerNA{}.
    \item[] Justification: \justification{While the scope of the paper is already very limited with regard to its experimental approach to Small Language Models, more details of the limitations are discussed in the Methods section.}
    \item[] Guidelines:
    \begin{itemize}
        \item The answer \answerNA{} means that the paper has no limitation while the answer \answerNo{} means that the paper has limitations, but those are not discussed in the paper. 
        \item The authors are encouraged to create a separate ``Limitations'' section in their paper.
        \item The paper should point out any strong assumptions and how robust the results are to violations of these assumptions (e.g., independence assumptions, noiseless settings, model well-specification, asymptotic approximations only holding locally). The authors should reflect on how these assumptions might be violated in practice and what the implications would be.
        \item The authors should reflect on the scope of the claims made, e.g., if the approach was only tested on a few datasets or with a few runs. In general, empirical results often depend on implicit assumptions, which should be articulated.
        \item The authors should reflect on the factors that influence the performance of the approach. For example, a facial recognition algorithm may perform poorly when image resolution is low or images are taken in low lighting. Or a speech-to-text system might not be used reliably to provide closed captions for online lectures because it fails to handle technical jargon.
        \item The authors should discuss the computational efficiency of the proposed algorithms and how they scale with dataset size.
        \item If applicable, the authors should discuss possible limitations of their approach to address problems of privacy and fairness.
        \item While the authors might fear that complete honesty about limitations might be used by reviewers as grounds for rejection, a worse outcome might be that reviewers discover limitations that aren't acknowledged in the paper. The authors should use their best judgment and recognize that individual actions in favor of transparency play an important role in developing norms that preserve the integrity of the community. Reviewers will be specifically instructed to not penalize honesty concerning limitations.
    \end{itemize}

\item {\bf Theory assumptions and proofs}
    \item[] Question: For each theoretical result, does the paper provide the full set of assumptions and a complete (and correct) proof?
    \item[] Answer: \answerNA{} % Replace by \answerYes{}, \answerNo{}, or \answerNA{}.
    \item[] Justification: \justification{This paper does not make use of any theoretical assumptions nor proofs.}
    \item[] Guidelines:
    \begin{itemize}
        \item The answer \answerNA{} means that the paper does not include theoretical results. 
        \item All the theorems, formulas, and proofs in the paper should be numbered and cross-referenced.
        \item All assumptions should be clearly stated or referenced in the statement of any theorems.
        \item The proofs can either appear in the main paper or the supplemental material, but if they appear in the supplemental material, the authors are encouraged to provide a short proof sketch to provide intuition. 
        \item Inversely, any informal proof provided in the core of the paper should be complemented by formal proofs provided in appendix or supplemental material.
        \item Theorems and Lemmas that the proof relies upon should be properly referenced. 
    \end{itemize}

    \item {\bf Experimental result reproducibility}
    \item[] Question: Does the paper fully disclose all the information needed to reproduce the main experimental results of the paper to the extent that it affects the main claims and/or conclusions of the paper (regardless of whether the code and data are provided or not)?
    \item[] Answer: \answerYes{} % Replace by \answerYes{}, \answerNo{}, or \answerNA{}.
    \item[] Justification: \justification{All of the code and data for training and prompting the models is supplied in the supplementary material. Additionally, all of the code and training data is freely accessible online under open licenses.}
    \item[] Guidelines:
    \begin{itemize}
        \item The answer \answerNA{} means that the paper does not include experiments.
        \item If the paper includes experiments, a \answerNo{} answer to this question will not be perceived well by the reviewers: Making the paper reproducible is important, regardless of whether the code and data are provided or not.
        \item If the contribution is a dataset and\slash or model, the authors should describe the steps taken to make their results reproducible or verifiable. 
        \item Depending on the contribution, reproducibility can be accomplished in various ways. For example, if the contribution is a novel architecture, describing the architecture fully might suffice, or if the contribution is a specific model and empirical evaluation, it may be necessary to either make it possible for others to replicate the model with the same dataset, or provide access to the model. In general. releasing code and data is often one good way to accomplish this, but reproducibility can also be provided via detailed instructions for how to replicate the results, access to a hosted model (e.g., in the case of a large language model), releasing of a model checkpoint, or other means that are appropriate to the research performed.
        \item While NeurIPS does not require releasing code, the conference does require all submissions to provide some reasonable avenue for reproducibility, which may depend on the nature of the contribution. For example
        \begin{enumerate}
            \item If the contribution is primarily a new algorithm, the paper should make it clear how to reproduce that algorithm.
            \item If the contribution is primarily a new model architecture, the paper should describe the architecture clearly and fully.
            \item If the contribution is a new model (e.g., a large language model), then there should either be a way to access this model for reproducing the results or a way to reproduce the model (e.g., with an open-source dataset or instructions for how to construct the dataset).
            \item We recognize that reproducibility may be tricky in some cases, in which case authors are welcome to describe the particular way they provide for reproducibility. In the case of closed-source models, it may be that access to the model is limited in some way (e.g., to registered users), but it should be possible for other researchers to have some path to reproducing or verifying the results.
        \end{enumerate}
    \end{itemize}


\item {\bf Open access to data and code}
    \item[] Question: Does the paper provide open access to the data and code, with sufficient instructions to faithfully reproduce the main experimental results, as described in supplemental material?
    \item[] Answer: \answerYes} % Replace by \answerYes{}, \answerNo{}, or \answerNA{}.
    \item[] Justification: \justification{All of the code for gathering and cleaning the training data, as well as hyperparameters for fine-tuning, and templates for prompting, is freely accessible online under open licenses.}
    \item[] Guidelines:
    \begin{itemize}
        \item The answer \answerNA{} means that paper does not include experiments requiring code.
        \item Please see the NeurIPS code and data submission guidelines (\url{https://neurips.cc/public/guides/CodeSubmissionPolicy}) for more details.
        \item While we encourage the release of code and data, we understand that this might not be possible, so \answerNo{} is an acceptable answer. Papers cannot be rejected simply for not including code, unless this is central to the contribution (e.g., for a new open-source benchmark).
        \item The instructions should contain the exact command and environment needed to run to reproduce the results. See the NeurIPS code and data submission guidelines (\url{https://neurips.cc/public/guides/CodeSubmissionPolicy}) for more details.
        \item The authors should provide instructions on data access and preparation, including how to access the raw data, preprocessed data, intermediate data, and generated data, etc.
        \item The authors should provide scripts to reproduce all experimental results for the new proposed method and baselines. If only a subset of experiments are reproducible, they should state which ones are omitted from the script and why.
        \item At submission time, to preserve anonymity, the authors should release anonymized versions (if applicable).
        \item Providing as much information as possible in supplemental material (appended to the paper) is recommended, but including URLs to data and code is permitted.
    \end{itemize}


\item {\bf Experimental setting/details}
    \item[] Question: Does the paper specify all the training and test details (e.g., data splits, hyperparameters, how they were chosen, type of optimizer) necessary to understand the results?
    \item[] Answer: \answerYes{} % Replace by \answerYes{}, \answerNo{}, or \answerNA{}.
    \item[] Justification: \justification{The body of the paper shares the most important information about the creation of the training dataset and some of the hyperparameters. The supplement supplies the full details of the training, and this information is also freely available online under open licenses.
    \item[] Guidelines:
    \begin{itemize}
        \item The answer \answerNA{} means that the paper does not include experiments.
        \item The experimental setting should be presented in the core of the paper to a level of detail that is necessary to appreciate the results and make sense of them.
        \item The full details can be provided either with the code, in appendix, or as supplemental material.
    \end{itemize}

\item {\bf Experiment statistical significance}
    \item[] Question: Does the paper report error bars suitably and correctly defined or other appropriate information about the statistical significance of the experiments?
    \item[] Answer: \answerNA{} % Replace by \answerYes{}, \answerNo{}, or \answerNA{}.
    \item[] Justification: \justification{This is a project about close reading analysis rather than model optimization. It is not a priority that the results of the experiment be optimizable or generalizable, so it does not report on statistical significance or error bars.}
    \item[] Guidelines:
    \begin{itemize}
        \item The answer \answerNA{} means that the paper does not include experiments.
        \item The authors should answer \answerYes{} if the results are accompanied by error bars, confidence intervals, or statistical significance tests, at least for the experiments that support the main claims of the paper.
        \item The factors of variability that the error bars are capturing should be clearly stated (for example, train/test split, initialization, random drawing of some parameter, or overall run with given experimental conditions).
        \item The method for calculating the error bars should be explained (closed form formula, call to a library function, bootstrap, etc.)
        \item The assumptions made should be given (e.g., Normally distributed errors).
        \item It should be clear whether the error bar is the standard deviation or the standard error of the mean.
        \item It is OK to report 1-sigma error bars, but one should state it. The authors should preferably report a 2-sigma error bar than state that they have a 96\% CI, if the hypothesis of Normality of errors is not verified.
        \item For asymmetric distributions, the authors should be careful not to show in tables or figures symmetric error bars that would yield results that are out of range (e.g., negative error rates).
        \item If error bars are reported in tables or plots, the authors should explain in the text how they were calculated and reference the corresponding figures or tables in the text.
    \end{itemize}

\item {\bf Experiments compute resources}
    \item[] Question: For each experiment, does the paper provide sufficient information on the computer resources (type of compute workers, memory, time of execution) needed to reproduce the experiments?
    \item[] Answer: \answerNA{} % Replace by \answerYes{}, \answerNo{}, or \answerNA{}.
    \item[] Justification: \justification{The language models used are very minor in scope (~1.5B in size), enough to run easily on ordinary consumer hardware. Therefore, the required computer resources are not reported in detail.}
    \item[] Guidelines:
    \begin{itemize}
        \item The answer \answerNA{} means that the paper does not include experiments.
        \item The paper should indicate the type of compute workers CPU or GPU, internal cluster, or cloud provider, including relevant memory and storage.
        \item The paper should provide the amount of compute required for each of the individual experimental runs as well as estimate the total compute. 
        \item The paper should disclose whether the full research project required more compute than the experiments reported in the paper (e.g., preliminary or failed experiments that didn't make it into the paper). 
    \end{itemize}
    
\item {\bf Code of ethics}
    \item[] Question: Does the research conducted in the paper conform, in every respect, with the NeurIPS Code of Ethics \url{https://neurips.cc/public/EthicsGuidelines}?
    \item[] Answer: \answerYes{} % Replace by \answerYes{}, \answerNo{}, or \answerNA{}.
    \item[] Justification: \justification{The research complies with all considerations surrounding privacy, security, fairness, safety, impact, etc.}
    \item[] Guidelines:
    \begin{itemize}
        \item The answer \answerNA{} means that the authors have not reviewed the NeurIPS Code of Ethics.
        \item If the authors answer \answerNo, they should explain the special circumstances that require a deviation from the Code of Ethics.
        \item The authors should make sure to preserve anonymity (e.g., if there is a special consideration due to laws or regulations in their jurisdiction).
    \end{itemize}


\item {\bf Broader impacts}
    \item[] Question: Does the paper discuss both potential positive societal impacts and negative societal impacts of the work performed?
    \item[] Answer: \answerYes{} % Replace by \answerYes{}, \answerNo{}, or \answerNA{}.
    \item[] Justification: \justification{The paper focuses on positive aspects of an overlooked aspect of language modeling. Although the paper mentions negative aspects related to large language modeling (i.e. the smoothing over of bias and the bigger-is-better mentality, both discussed in the introduction), it does not discuss these extensively.}
    \item[] Guidelines:
    \begin{itemize}
        \item The answer \answerNA{} means that there is no societal impact of the work performed.
        \item If the authors answer \answerNA{} or \answerNo, they should explain why their work has no societal impact or why the paper does not address societal impact.
        \item Examples of negative societal impacts include potential malicious or unintended uses (e.g., disinformation, generating fake profiles, surveillance), fairness considerations (e.g., deployment of technologies that could make decisions that unfairly impact specific groups), privacy considerations, and security considerations.
        \item The conference expects that many papers will be foundational research and not tied to particular applications, let alone deployments. However, if there is a direct path to any negative applications, the authors should point it out. For example, it is legitimate to point out that an improvement in the quality of generative models could be used to generate Deepfakes for disinformation. On the other hand, it is not needed to point out that a generic algorithm for optimizing neural networks could enable people to train models that generate Deepfakes faster.
        \item The authors should consider possible harms that could arise when the technology is being used as intended and functioning correctly, harms that could arise when the technology is being used as intended but gives incorrect results, and harms following from (intentional or unintentional) misuse of the technology.
        \item If there are negative societal impacts, the authors could also discuss possible mitigation strategies (e.g., gated release of models, providing defenses in addition to attacks, mechanisms for monitoring misuse, mechanisms to monitor how a system learns from feedback over time, improving the efficiency and accessibility of ML).
    \end{itemize}
    
\item {\bf Safeguards}
    \item[] Question: Does the paper describe safeguards that have been put in place for responsible release of data or models that have a high risk for misuse (e.g., pre-trained language models, image generators, or scraped datasets)?
    \item[] Answer: \answerNA{} % Replace by \answerYes{}, \answerNo{}, or \answerNA{}.
    \item[] Justification: \justification{The paper concerns very small and niche models that pose no (or extremely low) risks to safety.}
    \item[] Guidelines:
    \begin{itemize}
        \item The answer \answerNA{} means that the paper poses no such risks.
        \item Released models that have a high risk for misuse or dual-use should be released with necessary safeguards to allow for controlled use of the model, for example by requiring that users adhere to usage guidelines or restrictions to access the model or implementing safety filters. 
        \item Datasets that have been scraped from the Internet could pose safety risks. The authors should describe how they avoided releasing unsafe images.
        \item We recognize that providing effective safeguards is challenging, and many papers do not require this, but we encourage authors to take this into account and make a best faith effort.
    \end{itemize}

\item {\bf Licenses for existing assets}
    \item[] Question: Are the creators or original owners of assets (e.g., code, data, models), used in the paper, properly credited and are the license and terms of use explicitly mentioned and properly respected?
    \item[] Answer: \answerYes{} % Replace by \answerYes{}, \answerNo{}, or \answerNA{}.
    \item[] Justification: \justification{The author cites the sources for the dataset creation and the base model used for fine-tuning and respects copyright restrictions of the data.}
    \item[] Guidelines:
    \begin{itemize}
        \item The answer \answerNA{} means that the paper does not use existing assets.
        \item The authors should cite the original paper that produced the code package or dataset.
        \item The authors should state which version of the asset is used and, if possible, include a URL.
        \item The name of the license (e.g., CC-BY 4.0) should be included for each asset.
        \item For scraped data from a particular source (e.g., website), the copyright and terms of service of that source should be provided.
        \item If assets are released, the license, copyright information, and terms of use in the package should be provided. For popular datasets, \url{paperswithcode.com/datasets} has curated licenses for some datasets. Their licensing guide can help determine the license of a dataset.
        \item For existing datasets that are re-packaged, both the original license and the license of the derived asset (if it has changed) should be provided.
        \item If this information is not available online, the authors are encouraged to reach out to the asset's creators.
    \end{itemize}

\item {\bf New assets}
    \item[] Question: Are new assets introduced in the paper well documented and is the documentation provided alongside the assets?
    \item[] Answer: \answerYes{} % Replace by \answerYes{}, \answerNo{}, or \answerNA{}.
    \item[] Justification: \justification{The paper shares all assets freely online under open licenses.}
    \item[] Guidelines:
    \begin{itemize}
        \item The answer \answerNA{} means that the paper does not release new assets.
        \item Researchers should communicate the details of the dataset\slash code\slash model as part of their submissions via structured templates. This includes details about training, license, limitations, etc. 
        \item The paper should discuss whether and how consent was obtained from people whose asset is used.
        \item At submission time, remember to anonymize your assets (if applicable). You can either create an anonymized URL or include an anonymized zip file.
    \end{itemize}

\item {\bf Crowdsourcing and research with human subjects}
    \item[] Question: For crowdsourcing experiments and research with human subjects, does the paper include the full text of instructions given to participants and screenshots, if applicable, as well as details about compensation (if any)? 
    \item[] Answer: \answer{NA} % Replace by \answerYes{}, \answerNo{}, or \answerNA{}.
    \item[] Justification: \justification{The paper does not use crowdsourcing or human subjects.}
    \item[] Guidelines:
    \begin{itemize}
        \item The answer \answerNA{} means that the paper does not involve crowdsourcing nor research with human subjects.
        \item Including this information in the supplemental material is fine, but if the main contribution of the paper involves human subjects, then as much detail as possible should be included in the main paper. 
        \item According to the NeurIPS Code of Ethics, workers involved in data collection, curation, or other labor should be paid at least the minimum wage in the country of the data collector. 
    \end{itemize}

\item {\bf Institutional review board (IRB) approvals or equivalent for research with human subjects}
    \item[] Question: Does the paper describe potential risks incurred by study participants, whether such risks were disclosed to the subjects, and whether Institutional Review Board (IRB) approvals (or an equivalent approval/review based on the requirements of your country or institution) were obtained?
    \item[] Answer: \answer{NA} % Replace by \answerYes{}, \answerNo{}, or \answerNA{}.
    \item[] Justification: \justification{The paper does not involve crowdsourcing nor research with human subjects.}
    \item[] Guidelines:
    \begin{itemize}
        \item The answer \answerNA{} means that the paper does not involve crowdsourcing nor research with human subjects.
        \item Depending on the country in which research is conducted, IRB approval (or equivalent) may be required for any human subjects research. If you obtained IRB approval, you should clearly state this in the paper. 
        \item We recognize that the procedures for this may vary significantly between institutions and locations, and we expect authors to adhere to the NeurIPS Code of Ethics and the guidelines for their institution. 
        \item For initial submissions, do not include any information that would break anonymity (if applicable), such as the institution conducting the review.
    \end{itemize}

\item {\bf Declaration of LLM usage}
    \item[] Question: Does the paper describe the usage of LLMs if it is an important, original, or non-standard component of the core methods in this research? Note that if the LLM is used only for writing, editing, or formatting purposes and does \emph{not} impact the core methodology, scientific rigor, or originality of the research, declaration is not required.
    %this research? 
    \item[] Answer: \answerYes{} % Replace by \answerYes{}, \answerNo{}, or \answerNA{}.
    \item[] Justification: \justification{Prompting language models is a core part of the paper's methodology, which is described in detail throughout the paper. There is no other use of LLMs.}
    \item[] Guidelines:
    \begin{itemize}
        \item The answer \answerNA{} means that the core method development in this research does not involve LLMs as any important, original, or non-standard components.
        \item Please refer to our LLM policy in the NeurIPS handbook for what should or should not be described.
    \end{itemize}

\end{enumerate}
\end{document}
